Dans le circuit ci-dessous, les piles sont idéales : $\mathcal{E}_1 = \qty{\v{e1}}{V}$ et $\mathcal{E}_2 = \qty{\v{e2}}{V}$. Les résistances valent $R_1 = \qty{\v{R1}}{\ohm}$, $R_2 = \qty{\v{R2}}{\ohm}$ et $R_3 = \qty{\v{R3}}{\ohm}$. On note $I_1$ le courant qui monte dans la branche de gauche, $I_2$ le courant qui descend dans la branche du milieu et $I_3$ le courant qui monte dans la branche de droite.

\begin{center}
\begin{tikzpicture}[thick, scale=1.1]
    \draw (0,0) -- (0,0.9);
    \draw (-0.35,1.1) -- (0.35,1.1);
    \draw (-0.18,0.9) -- (0.18,0.9);
    \node at (-0.6,1.0) {$\mathcal{E}_1$};
    \node at (0.5,1.3) {\small $+$};
    \draw (0,1.1) -- (0,1.6);
    \draw[decorate,decoration={zigzag,segment length=5pt,amplitude=4pt}] (0,1.6) -- (0,2.6);
    \node at (-0.6,2.1) {$R_1$};
    \draw (0,2.6) -- (0,3.2) -- (6,3.2) -- (6,2.6);
    \draw[decorate,decoration={zigzag,segment length=5pt,amplitude=4pt}] (6,2.6) -- (6,1.6);
    \node at (6.6,2.1) {$R_3$};
    \draw (6,1.6) -- (6,1.1);
    \draw (5.65,1.1) -- (6.35,1.1);
    \draw (5.82,0.9) -- (6.18,0.9);
    \node at (6.6,1.0) {$\mathcal{E}_2$};
    \node at (5.5,1.3) {\small $+$};
    \draw (6,0.9) -- (6,0) -- (0,0);
    \draw (3,3.2) -- (3,2.1);
    \draw[decorate,decoration={zigzag,segment length=5pt,amplitude=4pt}] (3,2.1) -- (3,1.1);
    \node at (3.6,1.6) {$R_2$};
    \draw (3,1.1) -- (3,0);
    \fill (3,3.2) circle (2pt) node[above] {$A$};
    \fill (3,0) circle (2pt) node[below] {$B$};
    \draw[-{Stealth[length=3mm]}] (-0.4,2.9) -- (-0.4,3.4) node[left] {$I_1$};
    \draw[-{Stealth[length=3mm]}] (3.35,2.9) -- (3.35,2.4) node[right] {$I_2$};
    \draw[-{Stealth[length=3mm]}] (6.4,2.9) -- (6.4,3.4) node[right] {$I_3$};
\end{tikzpicture}
\end{center}

\begin{enumerate}[label=\alph*)]
    \item Écrivez la loi des nœuds au nœud $A$ et la loi des mailles pour les mailles de gauche et de droite.
    \item Déterminez le courant $I_2$ dans la résistance $R_2$.
    \item Déterminez les courants $I_1$ et $I_3$.
    \item Déterminez la puissance fournie par la pile $\mathcal{E}_1$.
\end{enumerate}
